\documentclass[11pt]{article}

\usepackage[a4paper,margin=1in]{geometry}
\usepackage{amsmath,amssymb,amsthm}
\usepackage[T1]{fontenc}
\usepackage{lmodern}
\usepackage{microtype}

\newcommand{\ph}{\varphi}
\newcommand{\N}{\mathbb N}
\newcommand{\Ssp}{\mathcal S}
\newcommand{\Ball}{B}
\newcommand{\Bell}[1]{B_{#1}}
\newcommand{\Delt}[1]{\Delta_{#1}}

\newtheorem*{theorem*}{Theorem}
\newtheorem*{lemma*}{Lemma}
\newtheorem*{proposition*}{Proposition}
\newtheorem*{remark*}{Remark}

\title{Self-Contained Better-Notation Rewrite of the Schlumprecht Szlenk-Radius Converse}
\author{Human-requested companion note}
\date{}

\begin{document}
\maketitle

\begin{center}
\fbox{\begin{minipage}{0.92\linewidth}
\textbf{Important note.}
This document is not the original LLM-generated output.  It is a
human-requested rewrite, requested by Antonio Acuaviva, of the proposed
argument with improved notation and with the details made self-contained.
The mathematical content is intended to be the same as in the original
solution packet, but the notation has been reorganized to make the proof
easier to check.
\end{minipage}}
\end{center}

\section{Setup}

Let
\[
 \ph(n)=\log_2(n+1)
 \qquad (n\in\N).
\]
Schlumprecht's space \(\Ssp\) is the completion of \(c_{00}\) under the norm
\[
 \|x\|_{\Ssp}
 =
 \max\left\{
 \|x\|_\infty,
 \sup_{n\ge 2,\ E_1<\cdots<E_n}
 \frac{1}{\ph(n)}
 \sum_{i=1}^n \|E_i x\|_{\Ssp}
 \right\}.
\]
Here \(E x\) denotes the coordinate restriction of \(x\) to the finite set
\(E\), and \(E_1<\cdots<E_n\) means the sets are successive.

We use the standard facts that the canonical basis of \(\Ssp\) is
1-unconditional and shrinking, and that \(\Ssp\) is reflexive.  Thus
\(e_N\to 0\) weakly in \(\Ssp\).

For a weak-star compact set \(K\subset X^*\), the Szlenk derivation is
\[
 s_\varepsilon(K)
 =
 \{x^*\in K:
 \operatorname{diam}(K\cap V)>\varepsilon
 \text{ for every weak-star neighbourhood }V\ni x^*\}.
\]
In the present case \(X=\Ssp^*\), so \(X^*=\Ssp\).  The corresponding
inner Szlenk radius is
\[
 r_{\Ssp^*}(\varepsilon)
 =
 \sup\{r>0:r\Ball_{\Ssp}\subset s_\varepsilon(\Ball_{\Ssp})\}.
\]

For \(0<\alpha<1\), define
\[
 a_\alpha
 =
 \inf_{n\in\N}
 \frac{\ph(n+1)-\alpha}{\ph(n)}.
\]
The known upper bound from the source paper is
\[
 r_{\Ssp^*}(\varepsilon)
 \le
 a_{\varepsilon/2}.
\]
The goal of the proposed argument is to prove the reverse inequality.

\begin{theorem*}
For \(0<\varepsilon<2\),
\[
 r_{\Ssp^*}(\varepsilon)
 =
 \inf_{n\in\N}
 \frac{\ph(n+1)-\varepsilon/2}{\ph(n)}.
\]
\end{theorem*}

\section{Logarithmic notation}

For \(i\ge 1\), write
\[
 \Bell{i}=\ph(i+1),
 \qquad
 \Delt{i}=\ph(i+1)-\ph(i).
\]
Thus
\[
 \Bell{i}=\log_2(i+2),
 \qquad
 \Delt{i}=\log_2\left(1+\frac{1}{i+1}\right).
\]
This notation is kept throughout the proof.  In particular, when the
product lemma is applied later with \(i=k-1\) and \(j=m\), the relevant
quantities remain visibly
\[
 \Delt{k-1},\qquad \Delt{m},\qquad \Bell{\ell},\qquad \Delt{\ell}.
\]

\section{The product lemma}

We first record the elementary monotonicity fact behind the product lemma.

\begin{lemma*}
The function
\[
 h(u)=\frac{u\log u}{u-1}
 \qquad (u>1)
\]
is increasing.
\end{lemma*}

\begin{proof}
Differentiate:
\[
 h'(u)
 =
 \frac{(u-1)(\log u+1)-u\log u}{(u-1)^2}
 =
 \frac{u-1-\log u}{(u-1)^2}.
\]
Since \(\log u\le u-1\), we have \(h'(u)\ge 0\).
\end{proof}

\begin{lemma*}[Product lemma]
For \(i,j\ge 1\), put
\[
 \ell=(i+1)(j+1)-1.
\]
Then
\[
 \Bell{\ell}\le \Bell{i}\Bell{j},
 \qquad
 \Delt{\ell}\le \Delt{i}\Delt{j}.
\]
\end{lemma*}

\begin{proof}
The first inequality is equivalent to
\[
 \log_2(pq+1)
 \le
 \log_2(p+1)\log_2(q+1)
 \qquad(p,q\ge 2),
\]
where \(p=i+1\) and \(q=j+1\).  For fixed \(p\), consider
\[
 F_p(q)=\frac{\log(pq+1)}{\log(q+1)}.
\]
A quotient-rule calculation gives
\[
 F_p'(q)
 =
 \frac{
 p(q+1)\log(q+1)-(pq+1)\log(pq+1)
 }{(pq+1)(q+1)(\log(q+1))^2}.
\]
The numerator is non-positive because it is the comparison
\[
 \frac{(q+1)\log(q+1)}{q}
 \le
 \frac{(pq+1)\log(pq+1)}{pq},
\]
which follows from the monotonicity of \(h\).  Hence \(F_p\) is decreasing,
so \(F_p(q)\le F_p(1)\).  This gives
\[
 \frac{\log(pq+1)}{\log(q+1)}
 \le
 \frac{\log(p+1)}{\log 2},
\]
which is the first inequality.

For the increment inequality, put \(a=1/p\) and \(b=1/q\).  It is enough to
prove
\[
 \log_2(1+ab)
 \le
 \log_2(1+a)\log_2(1+b)
 \qquad(0<a,b\le 1).
\]
For fixed \(a\), define
\[
 G_a(b)=\frac{\log(1+ab)}{\log(1+b)}.
\]
Again differentiating,
\[
 G_a'(b)
 =
 \frac{
 a(1+b)\log(1+b)-(1+ab)\log(1+ab)
 }{(1+ab)(1+b)(\log(1+b))^2}.
\]
Let \(u=1+b\) and \(v=1+ab\).  Since \(0<a\le 1\), \(v\le u\), and
\[
 a(u-1)=v-1.
\]
The numerator is
\[
 a u\log u-v\log v
 =
 (v-1)\left(
 \frac{u\log u}{u-1}
 -
 \frac{v\log v}{v-1}
 \right)\ge 0
\]
by the monotonicity of \(h\).  Hence \(G_a\) is increasing, so
\[
 G_a(b)\le G_a(1)=\frac{\log(1+a)}{\log 2}.
\]
This proves the increment inequality.
\end{proof}

\section{The scalar majorant}

First note that \(0<a_\alpha<1\).  For positivity, since
\(\ph(n+1)>\ph(n)\) and \(\ph(n)\ge 1\),
\[
 \frac{\ph(n+1)-\alpha}{\ph(n)}
 >
 \frac{\ph(n)-\alpha}{\ph(n)}
 =
 1-\frac{\alpha}{\ph(n)}
 \ge
 1-\alpha>0.
\]
Thus \(a_\alpha\ge 1-\alpha>0\).  For the upper bound, use
\(\ph(n+1)-\ph(n)\to 0\).  Choose \(n\) with
\[
 \ph(n+1)-\ph(n)<\alpha.
\]
Then
\[
 \frac{\ph(n+1)-\alpha}{\ph(n)}<1,
\]
so \(a_\alpha<1\).

From the definition of the infimum, for every \(n\in\N\),
\[
 a_\alpha
 \le
 \frac{\ph(n+1)-\alpha}{\ph(n)}.
\]
Equivalently,
\[
 \ph(n)a_\alpha+\alpha\le \ph(n+1).
\]

For \(u\ge 0\), define
\[
 M_\alpha(u)
 =
 \max\left\{
 u,\alpha,
 \sup_{m\in\N}
 \frac{\ph(m)u+\alpha}{\ph(m+1)}
 \right\}.
\]
This scalar function is designed to dominate the cost of adding one far
coordinate of size \(\alpha\).

If \(0\le u<a_\alpha\), then \(M_\alpha(u)\le 1\).  Indeed, \(u<1\) and
\(\alpha<1\), while for every \(m\),
\[
 \frac{\ph(m)u+\alpha}{\ph(m+1)}
 \le
 \frac{\ph(m)a_\alpha+\alpha}{\ph(m+1)}
 \le
 1.
\]

\section{Strengthened scalar estimate}

\begin{lemma*}[Strengthened scalar estimate]
Let \(0<\alpha<1\), \(0\le s\le t\), and \(k\ge 2\). Then
\[
 M_\alpha(s)
 +
 \min\{\ph(k-1)t,\ph(k)t-s\}
 \le
 \ph(k)M_\alpha(t).
\]
\end{lemma*}

\begin{proof}
We prove the estimate separately for each type of term appearing in
\(M_\alpha(s)\).

First consider the term \(s\).  Since
\[
 \min\{\ph(k-1)t,\ph(k)t-s\}\le \ph(k)t-s,
\]
we have
\[
 s+\min\{\ph(k-1)t,\ph(k)t-s\}
 \le
 \ph(k)t
 \le
 \ph(k)M_\alpha(t).
\]

Next consider the term \(\alpha\).  Since
\[
 \min\{\ph(k-1)t,\ph(k)t-s\}\le \ph(k-1)t,
\]
we get
\[
 \alpha+\min\{\ph(k-1)t,\ph(k)t-s\}
 \le
 \alpha+\ph(k-1)t.
\]
But \(M_\alpha(t)\) contains the affine term indexed by \(k-1\):
\[
 M_\alpha(t)
 \ge
 \frac{\ph(k-1)t+\alpha}{\ph(k)}.
\]
Thus
\[
 \alpha+\ph(k-1)t\le \ph(k)M_\alpha(t).
\]

It remains to consider an affine term indexed by \(m\):
\[
 \frac{\ph(m)s+\alpha}{\ph(m+1)}.
\]
For fixed \(t\), put
\[
 F_m(s)
 =
 \frac{\ph(m)s+\alpha}{\ph(m+1)}
 +
 \min\{\ph(k-1)t,\ph(k)t-s\}.
\]
The two quantities inside the minimum are equal when
\[
 \ph(k-1)t=\ph(k)t-s,
\]
that is, when
\[
 s=(\ph(k)-\ph(k-1))t=\Delt{k-1}t.
\]

On \(0\le s\le \Delt{k-1}t\), the minimum is \(\ph(k-1)t\), so \(F_m\) is
increasing.  On \(\Delt{k-1}t\le s\le t\), the minimum is \(\ph(k)t-s\),
and the slope of \(F_m\) is
\[
 \frac{\ph(m)}{\ph(m+1)}-1<0.
\]
Hence \(F_m\) is maximized at \(s=\Delt{k-1}t\).

At that point,
\[
 F_m(\Delt{k-1}t)
 =
 \frac{\ph(m)\Delt{k-1}t+\alpha}{\ph(m+1)}
 +
 \ph(k-1)t.
\]
Putting this over the denominator \(\ph(m+1)\), we get
\[
 F_m(\Delt{k-1}t)
 =
 \frac{
 \alpha
 +
 t\bigl(\ph(m)\Delt{k-1}
       +\ph(k-1)\ph(m+1)\bigr)
 }{\ph(m+1)}.
\]
Using
\[
 \ph(k)=\ph(k-1)+\Delt{k-1},
 \qquad
 \ph(m+1)=\ph(m)+\Delt{m},
\]
one obtains
\[
 \ph(m)\Delt{k-1}
 +\ph(k-1)\ph(m+1)
 =
 \ph(k)\ph(m+1)-\Delt{k-1}\Delt{m}.
\]
Therefore
\[
 F_m(\Delt{k-1}t)
 =
 \frac{
 \alpha
 +
 t\bigl(\ph(k)\ph(m+1)-\Delt{k-1}\Delt{m}\bigr)
 }{\ph(m+1)}.
\]
It is enough to prove
\[
 M_\alpha(t)
 \ge
 t+
 \frac{\alpha-\Delt{k-1}\Delt{m}t}
      {\ph(k)\ph(m+1)}.
\]

Apply the product lemma with \(i=k-1\) and \(j=m\).  Then
\[
 \ell=(k-1+1)(m+1)-1=k(m+1)-1,
\]
and
\[
 \Bell{\ell}\le \Bell{k-1}\Bell{m}
 =
 \ph(k)\ph(m+1),
\]
while
\[
 \Delt{\ell}\le \Delt{k-1}\Delt{m}.
\]
The affine term indexed by \(\ell\) in \(M_\alpha(t)\) gives
\[
 M_\alpha(t)
 \ge
 \frac{\ph(\ell)t+\alpha}{\ph(\ell+1)}.
\]
Since
\[
 \Bell{\ell}=\ph(\ell+1),
 \qquad
 \Delt{\ell}=\ph(\ell+1)-\ph(\ell),
\]
we have
\[
 \ph(\ell)=\Bell{\ell}-\Delt{\ell}.
\]
Hence
\[
 M_\alpha(t)
 \ge
 t+\frac{\alpha-\Delt{\ell}t}{\Bell{\ell}}.
\]

If
\[
 \alpha-\Delt{k-1}\Delt{m}t<0,
\]
then
\[
 t+
 \frac{\alpha-\Delt{k-1}\Delt{m}t}
      {\ph(k)\ph(m+1)}
 <t,
\]
and the desired inequality follows from \(M_\alpha(t)\ge t\).

If
\[
 \alpha-\Delt{k-1}\Delt{m}t\ge 0,
\]
then
\[
 \alpha-\Delt{\ell}t
 \ge
 \alpha-\Delt{k-1}\Delt{m}t
 \ge 0,
\]
and
\[
 \Bell{\ell}\le \ph(k)\ph(m+1).
\]
Therefore
\[
 \frac{\alpha-\Delt{\ell}t}{\Bell{\ell}}
 \ge
 \frac{\alpha-\Delt{k-1}\Delt{m}t}
      {\ph(k)\ph(m+1)}.
\]
Thus the affine-term case follows.

The estimate has now been checked for \(s\), for \(\alpha\), and for every
affine term indexed by \(m\).  Hence it holds for \(M_\alpha(s)\).
\end{proof}

\section{A comparison lemma for the Schlumprecht norm}

Define iterative norms on \(c_{00}\) by
\[
 \|x\|_0=\|x\|_\infty,
\]
and
\[
 \|x\|_{r+1}
 =
 \max\left\{
 \|x\|_\infty,
 \sup_{n\ge 2,\ E_1<\cdots<E_n}
 \frac{1}{\ph(n)}
 \sum_{i=1}^n \|E_i x\|_r
 \right\}.
\]
Then
\[
 \|x\|_{\Ssp}=\sup_{r\ge 0}\|x\|_r.
\]

\begin{lemma*}[Local comparison]
Let \(y\in c_{00}\).  Suppose there is a function
\[
 \rho:[\N]^{<\omega}\to [0,\infty),
\]
where
\[
 [\N]^{<\omega}=\{F\subset \N:F\text{ is finite}\},
\]
such that, for every finite \(F\subset\N\),
\[
 \|Fy\|_\infty\le \rho(F),
\]
for every \(G\subset F\),
\[
 G\subset F \Longrightarrow \rho(G)\le \rho(F),
\]
and, whenever \(k\ge 2\) and \(E_1<\cdots<E_k\subset F\),
\[
 \sum_{i=1}^k \rho(E_i)\le \ph(k)\rho(F).
\]
Then
\[
 \|Fy\|_{\Ssp}\le \rho(F)
\]
for every finite \(F\subset\N\).
\end{lemma*}

\begin{proof}
We prove by induction on \(r\) that
\[
 \|Fy\|_r\le \rho(F)
\]
for every finite \(F\).  The case \(r=0\) is exactly the assumption
\(\|Fy\|_\infty\le \rho(F)\).

Assume the claim holds for \(r\).  For \(E_1<\cdots<E_k\subset F\),
\[
 \sum_{i=1}^k \|E_i y\|_r
 \le
 \sum_{i=1}^k \rho(E_i)
 \le
 \ph(k)\rho(F).
\]
Together with \(\|Fy\|_\infty\le\rho(F)\), this gives
\[
 \|Fy\|_{r+1}\le \rho(F).
\]
Taking the supremum over \(r\) proves the result.
\end{proof}

\section{Far-coordinate perturbation}

\begin{proposition*}
Let \(0<\alpha<1\), let \(x\in c_{00}\), and assume
\[
 \|x\|_{\Ssp}<a_\alpha.
\]
If \(N>\max\operatorname{supp}x\), then
\[
 \|x+\alpha e_N\|_{\Ssp}\le 1
 \qquad\text{and}\qquad
 \|x-\alpha e_N\|_{\Ssp}\le 1.
\]
\end{proposition*}

\begin{proof}
By 1-unconditionality, it is enough to prove the estimate for
\(|x|+\alpha e_N\).  Thus assume \(x\ge 0\), and set
\[
 y=x+\alpha e_N.
\]

For each finite \(F\subset\N\), define
\[
 \rho(F)
 =
 \begin{cases}
 \|Fx\|_{\Ssp},& N\notin F,\\[2mm]
 M_\alpha(\|Fx\|_{\Ssp}),& N\in F.
 \end{cases}
\]
The lattice monotonicity of the Schlumprecht norm and the monotonicity of
\(M_\alpha\) give
\[
 G\subset F\Longrightarrow \rho(G)\le \rho(F).
\]
Also,
\[
 \|Fy\|_\infty\le \rho(F),
\]
because \(M_\alpha(u)\ge u\) and \(M_\alpha(u)\ge\alpha\).

It remains to verify the local comparison inequality.  Fix finite
\(F\subset\N\), and let
\[
 E_1<\cdots<E_k\subset F,
 \qquad k\ge 2.
\]

If \(N\notin F\), then
\[
 \sum_{i=1}^k \rho(E_i)
 =
 \sum_{i=1}^k \|E_i x\|_{\Ssp}
 \le
 \ph(k)\|Fx\|_{\Ssp}
 =
 \ph(k)\rho(F).
\]

Assume now that \(N\in F\).
If none of the \(E_i\) contains \(N\), then
\[
 \sum_{i=1}^k \rho(E_i)
 =
 \sum_{i=1}^k \|E_i x\|_{\Ssp}
 \le
 \ph(k)\|Fx\|_{\Ssp}
 \le
 \ph(k)M_\alpha(\|Fx\|_{\Ssp})
 =
 \ph(k)\rho(F).
\]

Finally suppose exactly one set \(E_j\) contains \(N\), and fix this \(j\).
Then
\[
 \sum_{i=1}^k \rho(E_i)
 =
 M_\alpha(\|E_jx\|_{\Ssp})+\sum_{i\ne j}\|E_i x\|_{\Ssp}.
\]
The remaining pieces satisfy two estimates.  First,
\[
 \sum_{i\ne j}\|E_i x\|_{\Ssp}\le \ph(k-1)\|Fx\|_{\Ssp}.
\]
For \(k=2\), this is just monotonicity, since \(\ph(1)=1\).  For
\(k>2\), it follows from the Schlumprecht recursive inequality applied to
the \(k-1\) remaining blocks.  Second, using the full \(k\)-tuple,
\[
 \|E_jx\|_{\Ssp}+\sum_{i\ne j}\|E_i x\|_{\Ssp}
 \le
 \ph(k)\|Fx\|_{\Ssp},
\]
so
\[
 \sum_{i\ne j}\|E_i x\|_{\Ssp}
 \le
 \ph(k)\|Fx\|_{\Ssp}-\|E_jx\|_{\Ssp}.
\]
Therefore
\[
 \sum_{i\ne j}\|E_i x\|_{\Ssp}
 \le
 \min\{\ph(k-1)\|Fx\|_{\Ssp},
          \ph(k)\|Fx\|_{\Ssp}-\|E_jx\|_{\Ssp}\}.
\]
By the strengthened scalar estimate,
\[
\begin{aligned}
 \sum_{i=1}^k \rho(E_i)
 &\le
 M_\alpha(\|E_jx\|_{\Ssp})+
 \min\{\ph(k-1)\|Fx\|_{\Ssp},
        \ph(k)\|Fx\|_{\Ssp}-\|E_jx\|_{\Ssp}\} \\
 &\le
 \ph(k)M_\alpha(\|Fx\|_{\Ssp})
 =
 \ph(k)\rho(F).
\end{aligned}
\]

The local comparison lemma gives
\[
 \|Fy\|_{\Ssp}\le \rho(F)
\]
for every finite \(F\).  Taking \(F=\operatorname{supp}x\cup\{N\}\), we get
\[
 \|x+\alpha e_N\|_{\Ssp}
 \le
 M_\alpha(\|x\|_{\Ssp})
 \le 1,
\]
because \(\|x\|_{\Ssp}<a_\alpha\).  The estimate for
\(x-\alpha e_N\) follows by unconditionality.
\end{proof}

\section{The reverse Szlenk-radius inclusion}

Let \(0<\varepsilon<2\), and put
\[
 \beta=\varepsilon/2.
\]
We prove
\[
 r_{\Ssp^*}(\varepsilon)\ge a_\beta.
\]

First take \(x\in c_{00}\) with
\[
 \|x\|_{\Ssp}<a_\beta.
\]
Choose \(\alpha\) such that
\[
 \beta<\alpha<1
 \qquad\text{and}\qquad
 \|x\|_{\Ssp}<a_\alpha.
\]
This is possible because the map \(\alpha\mapsto a_\alpha\) is continuous:
it is the infimum of affine functions
\[
 \alpha\mapsto \frac{\ph(n+1)-\alpha}{\ph(n)}
\]
whose slopes have absolute value at most \(1\).

For every sufficiently large \(N\), the far-coordinate perturbation result
gives
\[
 x+\alpha e_N\in \Ball_{\Ssp},
 \qquad
 x-\alpha e_N\in \Ball_{\Ssp}.
\]
Since \(e_N\to 0\) weakly, both points converge weakly to \(x\).  Their
distance is
\[
 \|(x+\alpha e_N)-(x-\alpha e_N)\|_{\Ssp}
 =
 \|2\alpha e_N\|_{\Ssp}
 =
 2\alpha
 >
 \varepsilon.
\]
Therefore every weak neighbourhood of \(x\) meets \(\Ball_{\Ssp}\) in a set
of diameter \(>\varepsilon\).  Hence
\[
 x\in s_\varepsilon(\Ball_{\Ssp}).
\]

Now let \(r<a_\beta\).  If \(z\in r\Ball_{\Ssp}\), choose finitely
supported \(x_n\to z\) in norm with \(\|x_n\|_{\Ssp}<a_\beta\) for all
large \(n\).  By the finite-support case,
\[
 x_n\in s_\varepsilon(\Ball_{\Ssp})
\]
eventually.  The set \(s_\varepsilon(\Ball_{\Ssp})\) is weakly closed,
hence norm closed, so
\[
 z\in s_\varepsilon(\Ball_{\Ssp}).
\]
Thus
\[
 r\Ball_{\Ssp}\subset s_\varepsilon(\Ball_{\Ssp})
 \qquad(r<a_\beta),
\]
and therefore
\[
 r_{\Ssp^*}(\varepsilon)\ge a_\beta.
\]
Together with the known upper bound, this gives
\[
 r_{\Ssp^*}(\varepsilon)
 =
 a_{\varepsilon/2}
 =
 \inf_{n\in\N}
 \frac{\ph(n+1)-\varepsilon/2}{\ph(n)}.
\]

\section*{Comment on notation}

The original proof introduced several temporary letters for quantities such
as \(\ph(k)\), \(\ph(m+1)\), \(\Delt{k-1}\), \(\Delt{m}\),
\(\Bell{\ell}\), and \(\Delt{\ell}\).  This rewrite keeps the logarithmic
weights and increments visible.  In particular, the key use of the product
lemma in the scalar estimate is simply
\[
 \Bell{\ell}\le \ph(k)\ph(m+1),
 \qquad
 \Delt{\ell}\le \Delt{k-1}\Delt{m},
\]
where
\[
 \ell=k(m+1)-1.
\]

\end{document}
